% BHU PYQ -- Transcribed & Corrected
% Paper: SOCMD31: Cinema & Society
% Exam: Under Graduate Semester - III Examination, 2025-26 (NEP)
% Ref Code: D/180
% Category: Multidisciplinary Course (MDC)
% Department: Sociology / Faculty of Social Sciences

\documentclass[11pt,a4paper]{article}
\usepackage[a4paper,top=2.5cm,bottom=2.5cm,left=2.8cm,right=2.8cm,headheight=14pt]{geometry}
\usepackage{amsmath,amssymb}
\usepackage{fontspec}
\setmainfont{Kohinoor Devanagari}
\usepackage{microtype,fancyhdr,enumitem,hyperref,lastpage,parskip}

\hypersetup{
    colorlinks=true,
    urlcolor=black,
    linkcolor=black,
    pdftitle={SOCMD31: Cinema & Society -- BHU Sem III 2025-26 (NEP MDC)}
}

\pagestyle{fancy}
\fancyhf{}
\renewcommand{\headrulewidth}{0.4pt}
\renewcommand{\footrulewidth}{0.4pt}
\fancyhead[L]{\small \textit{Banaras Hindu University}}
\fancyhead[R]{\small \textit{SOCMD31: Cinema \& Society (2025-26)}}
\fancyfoot[C]{\small Page \thepage\ of \pageref{LastPage}}

\newcommand{\pts}[1]{\hfill \textit{[#1]}}

\begin{document}

\begin{center}
  {\large\bfseries BANARAS HINDU UNIVERSITY}\\[2pt]
  {\normalsize Under Graduate Semester III Examination, 2025-26 (NEP)}\\[6pt]
  \rule{0.8\linewidth}{0.5pt}\\[6pt]
  {\large\bfseries SOCIOLOGY (MDC)}\\[4pt]
  {\large\bfseries Paper Code: SOCMD31 : Cinema \& Society (सिनेमा और समाज)}\\[2pt]
  {\small Ref. No.: D/180}\\[4pt]
  {\normalsize Time: 2$\frac{1}{2}$ Hours\hfill Maximum Marks: 70}
\end{center}

\rule{\linewidth}{0.4pt}

\smallskip

\noindent \textit{Note: This question-paper consists of three Sections: A, B and C. Answer the questions as per the instructions given in each section. (इस प्रश्न-पत्र में तीन खण्ड : अ, ब और स हैं। प्रत्येक खण्ड में दिये गये निर्देशानुसार प्रश्नों के उत्तर दीजिए।)}

\bigskip

\begin{center}
  \textbf{SECTION -- A / खण्ड -- अ}\\
  \textit{(Objective Type Questions / वस्तुनिष्ठ प्रश्न)}
\end{center}
\noindent \textit{Answer each within 50 words. Each question carries 2 marks. ($2 \times 5 = 10$ marks)}

\begin{enumerate}[label=\textbf{\arabic*.}]
  \item Answer the following (निम्नलिखित के उत्तर दीजिए) :
  \begin{enumerate}[label=(\roman*)]
    \item Mention title of any two movies that promote sports culture in India. (भारत में खेल संस्कृति को बढ़ावा देने वाली किन्हीं दो फिल्मों के शीर्षक बताइए।)
    \item Who is the author of the book titled ``What is Cinema?'' published in 2004 by University of California Press? (कैलिफोर्निया विश्वविद्यालय प्रेस द्वारा 2004 में प्रकाशित ``सिनेमा क्या है?'' नामक पुस्तक के लेखक कौन हैं?)
    \item Write title of any two silent movies. (किन्हीं दो मूक फिल्मों के शीर्षक लिखिए।)
    \item Write names of any two Indian cinemas that attracted global audience. (विश्व भर के दर्शकों को आकर्षित करने वाले किन्हीं दो भारतीय सिनेमा के नाम लिखिए।)
    \item Mention any two famous international film awards. (किन्हीं दो प्रसिद्ध अंतर्राष्ट्रीय फिल्म पुरस्कारों का उल्लेख कीजिए।)
  \end{enumerate}
\end{enumerate}

\bigskip

\begin{center}
  \textbf{SECTION -- B / खण्ड -- ब}\\
  \textit{(Short Answer Type Questions / अतिलघु उत्तरीय प्रश्न)}
\end{center}
\noindent \textit{Answer each within 250 words. Each question carries 10 marks. ($10 \times 3 = 30$ marks)}

\begin{enumerate}[label=\textbf{\arabic*.}, start=2]
  \item Write a note on silent films. (मूक फिल्मों पर एक टिप्पणी लिखिए।)\pts{10}\\
  \begin{center}\textbf{OR / अथवा}\end{center}
  Explain the relationship between cinema and society with suitable examples. (सिनेमा और समाज के बीच संबंधों को उपयुक्त उदाहरणों के साथ समझाइये।)

  \item ``Cinema is considered as a vehicle of ideology/ies.'' Justify the statement. (``सिनेमा को विचारधारा/विचारधाराओं का वाहन माना जाता है।'' इस कथन का औचित्य सिद्ध कीजिए।)\pts{10}\\
  \begin{center}\textbf{OR / अथवा}\end{center}
  Describe the representation of female characters and their roles in Indian cinemas. (भारतीय सिनेमा में महिला पात्रों के प्रतिनिधित्व और उनकी भूमिकाओं का वर्णन करें।)

  \item Explain the impact of digital technology on cinema. (सिनेमा पर डिजिटल प्रौद्योगिकी के प्रभाव की व्याख्या करें।)\pts{10}\\
  \begin{center}\textbf{OR / अथवा}\end{center}
  Describe the factors responsible for global spread of cinema. (सिनेमा के वैश्विक प्रसार के लिए उत्तरदायी कारकों का वर्णन कीजिए।)
\end{enumerate}

\bigskip

\begin{center}
  \textbf{SECTION -- C / खण्ड -- स}\\
  \textit{(Long Answer Type Questions / दीर्घ उत्तरीय प्रश्न)}
\end{center}
\noindent \textit{Answer all the questions. Each question carries 15 marks. ($15 \times 2 = 30$ marks)}

\begin{enumerate}[label=\textbf{\arabic*.}, start=5]
  \item ``The democratization of film production, driven by affordable technology and digital platforms, has significant positive and negative impacts on society.'' Substantiate the statement. (``सस्ती तकनीक और डिजिटल प्लेटफॉर्म द्वारा संचालित फिल्म निर्माण के लोकतंत्रीकरण का समाज पर महत्त्वपूर्ण सकारात्मक और नकारात्मक प्रभाव पड़ता है।'' इस कथन की पुष्टि कीजिए।)\pts{15}\\
  \begin{center}\textbf{OR / अथवा}\end{center}
  What do you mean by cinema? Write an essay on the birth of cinema. (सिनेमा से आप क्या समझते हैं? सिनेमा के जन्म पर एक निबन्ध लिखिए।)

  \item Discuss how piracy may offer minor benefits to viewers, yet cause significant harm to the film industry. Also suggest five measures to reduce the 'culture of piracy'. (चर्चा कीजिए कि कैसे पायरेसी दर्शकों को मामूली लाभ पहुँचाती है, जबकि फिल्म उद्योग को भारी नुकसान पहुँचाती है। साथ ही, 'पायरेसी की संस्कृति' को कम करने के पाँच उपाय भी सुझाइए।)\pts{15}\\
  \begin{center}\textbf{OR / अथवा}\end{center}
  Distinguish between Bollywood movies and other regional movies of India from sociological and cultural perspectives. (समाजशास्त्रीय और सांस्कृतिक दृष्टिकोण से बॉलीवुड फिल्मों और भारत की अन्य क्षेत्रीय फिल्मों के बीच अंतर स्पष्ट करें।)
\end{enumerate}

\bigskip
\begin{center}
  \textbf{***}
\end{center}

\end{document}
